\chapter*{Dossier grand oral}
\subsection*{\centering Objectif : Comprendre l'épreuve.}

L'épreuve du Grand oral permet au candidat de montrer sa capacité à prendre la parole en public de façon claire et convaincante. Elle lui offre aussi l’opportunité d'utiliser les connaissances liées à ses spécialités pour démontrer ses capacités argumentatives.\\

Le candidat présente au jury \textbf{deux} questions préparées avec ses professeurs qui portent sur ses deux spécialités, soit prises isolément, soit abordées de manière transversale en voie générale.\\

Le jury choisit une de ces deux questions. Le candidat a ensuite 20 minutes de préparation pour mettre en ordre ses idées et préparer s’il le souhaite un support sur du papier. Ce support est une aide pour la parole du candidat ; il n’a pas vocation à être donné à lire au jury.
Il s’agit de notes, d’un plan d’exposé, de trame de prise de parole, de mots-clefs ou d’idées directrices.
Ces notes peuvent aussi servir de document d’appui à l’argumentation (schéma, courbe, diagramme, tableau, formule mathématique…).

\section{Modalités}
\begin{description}
    \item[Temps de préparation] : 20 min
    \item[Durée de l'épreuve] : 20 min, soit 5 min de présentation de la question (sans être interrompu·e), 10 min d'échange avec le jury sur cette question, puis 5 min d'échange sur le projet d'orientation
    \item[Support]: Pas de support visuel, seule une feuille de note avec votre plan que vous aurez rédigé durant le temps de préparation est autorisé. 
    \item[Coefficient]: 10 (soit 200 pts sur 2000 possibles, le deuxième plus fort coefficient derrière les spécialités).
    \item[Jury]: Le jury est composé de deux professeurs de disciplines différentes, dont l'un représente l'un des deux enseignements de spécialité du candidat et l'autre représente l'autre enseignement de spécialité ou l'un des enseignements communs, ou est professeur·e-documentaliste 
\end{description}



Les ressources officielles :
\begin{itemize}
    \item \url{https://eduscol.education.fr/document/7136/download}
    \item \url{https://eduscol.education.fr/media/3919/download}
\end{itemize}

\section{Évaluation}
Le jury évalue l'épreuve avec la \textbf{grille officielle d'évaluation du Grand oral}, publiée par le ministère : elle note l'épreuve sur 20 points, selon plusieurs critères. Lisez la grille complète sur la page d'Éduscol :\\
\url{https://eduscol.education.fr/729/presentation-du-grand-oral}
\newpage





